% Optional formatting reference; not included in main.tex.
\documentclass[11pt]{article}
\usepackage{problemset}
% Edit these settings for each assignment.
\newcommand{\authorname}{Generic}
\newcommand{\problemsetnumber}{1}
\newcommand{\assignmentdate}{\today} % Automatically uses the compilation date.


\begin{document}
\begin{problem}[Lists and mathematics]\label{prob:example}
Let \(x\in\R\).
\begin{enumerate}
  \item Verify the following identity.
  \begin{enumerate}
    \item Expand the square.
    \item Collect terms.
  \end{enumerate}
  \item Explain why \(x^2\geq 0\).
\end{enumerate}
\[
  (x+1)^2=x^2+2x+1.
\]
\end{problem}

\begin{solution}
\begin{proof}
By the distributive law,
\begin{align}
  (x+1)^2 &= (x+1)(x+1) \\
          &= x^2+2x+1.
\end{align}
If \(x\geq0\), the product \(x\cdot x\) is nonnegative.
If \(x<0\), then \(-x>0\) and \(x^2=(-x)^2>0\).
\end{proof}
\end{solution}

\begin{problem}[Long statements and page breaks]
This deliberately long placeholder illustrates a statement that can continue
across pages. Replace it with your own question. The box follows the text
automatically; no fixed height or manual page break is needed.

\newcount\exampleparagraph
\exampleparagraph=0
\loop
  \advance\exampleparagraph by 1
  This is placeholder paragraph \the\exampleparagraph. A longer question may
  introduce notation, list assumptions, describe a procedure, and specify
  what the solution should establish. Keep these instructions in ordinary
  paragraphs so the statement can break naturally across pages. Separate
  assumptions from the requested conclusion and define symbols before using
  them. This text demonstrates layout only and contains no assignment content.
  \par
\ifnum\exampleparagraph<9
\repeat
\end{problem}

\begin{solution}
Write your explanation here. Refer to earlier questions with
\verb|\ref|, for example Problem~\ref{prob:example}.

\begin{figure}[htbp]
  \centering
  \fbox{\parbox[c][1in][c]{0.75\linewidth}{\centering Figure placeholder}}
  \caption{Replace the box with your own graphic using \texttt{\string\includegraphics}.}
\end{figure}

\begin{table}[htbp]
  \centering
  \begin{tabular}{@{}ll@{}}
    \toprule
    Item & Description \\
    \midrule
    Placeholder & Replace with your own content \\
    \bottomrule
  \end{tabular}
  \caption{Table placeholder; no numerical results are supplied.}
\end{table}
\end{solution}
\end{document}
